% TODO checklist:
% - [x] README.md: UI descriptions, technology stack
% - [x] Q&A.md (Section H, Q47-Q49): UI architectures, API interaction, dual-UI analysis
% - [x] ui_agent/: Next.js application structure, components, state management
% - [x] ui_control_panel/: Streamlit application structure, views, API client

\section{Presentation Layer}
\label{sec:presentation-layer}

The Presentation Layer provides two specialized user interfaces optimized for distinct personas and use cases: a \textbf{Next.js Agent Chat UI} for end-user conversational interactions, and a \textbf{Streamlit Control Panel} for administrative operations and system monitoring. This dual-UI strategy enables persona-specific optimization while sharing the common backend API.

\subsection{Agent Chat UI}
\label{subsec:agent-chat-ui}

The Agent Chat UI is a modern, production-grade web application built with Next.js 14, providing a Copilot-style conversational interface for interacting with AI agents.

\paragraph{Technology Stack}
Table~\ref{tab:chat-ui-stack} summarizes the technologies used in the Agent Chat UI.

\begin{table}[htbp]
    \centering
    \caption{Agent Chat UI technology stack.}
    \label{tab:chat-ui-stack}
    \small
    \begin{tabular}{lll}
        \toprule
        \textbf{Component} & \textbf{Technology} & \textbf{Version} \\
        \midrule
        Framework & Next.js (App Router) & 14.2.15 \\
        UI Library & React & 18.3.1 \\
        State Management & Zustand & 4.5.5 \\
        UI Components & Radix UI & Latest \\
        Styling & Tailwind CSS & 3.4.14 \\
        Icons & Lucide React & Latest \\
        Type Safety & TypeScript & 5.x \\
        \bottomrule
    \end{tabular}
\end{table}

\paragraph{Component Architecture}
The UI follows a component-based architecture with clear separation of concerns:

\begin{lstlisting}[language=bash, caption={Agent Chat UI directory structure.}]
ui_agent/src/
|-- app/
|   |-- page.tsx          # Main chat page with layout
|   |-- layout.tsx        # Root layout with providers
|   |-- globals.css       # Global styles (Tailwind)
|   `-- api/              # API routes (auth token proxy)
|-- components/
|   |-- chat-area.tsx     # Message display + step rendering
|   |-- chat-input.tsx    # Prompt input + model selector
|   |-- role-toggle.tsx   # Admin/User role switch
|   `-- ui/               # Radix UI primitives
|-- stores/
|   |-- auth-store.ts     # Authentication state (Zustand + persist)
|   `-- chat-store.ts     # Chat messages and runs (Zustand)
`-- lib/
    |-- api.ts            # Backend API client
    `-- utils.ts          # Utility functions
\end{lstlisting}

\paragraph{State Management}
The application uses Zustand for lightweight, performant state management with two primary stores:

\begin{description}
    \item[Auth Store] Manages role selection (Admin/User), token generation via Auth0, and SSR-safe hydration. Role selection persists in localStorage, but tokens are regenerated on each session for security.
    
    \item[Chat Store] Manages chat messages, agent runs with steps and status, model selection, and auto-scroll behavior. Messages are stored as a flat array with run references for efficient rendering.
\end{description}

\begin{lstlisting}[language=JavaScript, caption={Zustand store interface pattern.}]
interface ChatState {
  messages: ChatMessage[];
  selectedModel: string;
  availableModels: Array<{ id: string; name: string }>;
  currentRunId: string | null;
  isSubmitting: boolean;
  isPolling: boolean;
  
  addUserMessage: (content: string) => string;
  addAgentResponse: (messageId: string, run: AgentRun) => void;
  updateAgentRun: (messageId: string, run: AgentRun) => void;
}
\end{lstlisting}

\paragraph{Backend API Interaction}
The API client in \texttt{lib/api.ts} provides typed functions for all backend communication:

\begin{table}[htbp]
    \centering
    \caption{Agent Chat UI API endpoints.}
    \label{tab:chat-api-endpoints}
    \small
    \begin{tabular}{lll}
        \toprule
        \textbf{Endpoint} & \textbf{Function} & \textbf{Purpose} \\
        \midrule
        \texttt{POST /v1/agent-runs} & \texttt{createAgentRun()} & Create new agent run \\
        \texttt{GET /v1/agent-runs/\{id\}} & \texttt{getAgentRun()} & Poll run status \\
        \texttt{GET /v1/agent-runs/\{id\}/steps} & \texttt{getAgentRunSteps()} & Get execution steps \\
        \texttt{GET /v1/models/instances} & \texttt{listModels()} & List available models \\
        \texttt{GET /v1/models/defaults} & \texttt{getDefaultModel()} & Get default model \\
        \texttt{GET /v1/auth/me} & \texttt{getAuthMe()} & Validate token \\
        \bottomrule
    \end{tabular}
\end{table}

\paragraph{Polling Pattern}
The chat UI implements a time-based polling pattern for agent run completion:

\begin{lstlisting}[language=JavaScript, caption={Run completion polling implementation.}]
async function pollRunUntilComplete(
  runId: string,
  token: string | null,
  onUpdate: (run: AgentRun) => void,
  intervalMs: number = 2000,    // 2 second intervals
  maxAttempts: number = 300     // 10 minutes max
): Promise<AgentRun> {
  let attempts = 0;
  
  while (attempts < maxAttempts) {
    const run = await getAgentRun(runId, token);
    onUpdate(run);  // Update UI with latest state
    
    // Terminal states
    if (['succeeded', 'failed', 'cancelled'].includes(run.status)) {
      return run;
    }
    
    await new Promise(resolve => setTimeout(resolve, intervalMs));
    attempts++;
  }
  
  throw new Error(`Polling timed out after ${maxAttempts} attempts`);
}
\end{lstlisting}

\paragraph{Model Selection Flow}
The chat input component handles dynamic model selection:
\begin{enumerate}
    \item On role change, fetch available models from \texttt{/v1/models/instances}
    \item Query \texttt{/v1/models/defaults} for backend-configured default
    \item If backend default exists, use it; otherwise use first available model
    \item User can override via dropdown selection
\end{enumerate}

\paragraph{Message Display Features}
The chat area component renders messages with rich step visualization:
\begin{itemize}
    \item \textbf{User messages}: Simple text bubbles with timestamp
    \item \textbf{Agent responses}: Full run timeline including:
    \begin{itemize}
        \item Status indicators (loading spinner, success checkmark, error icon)
        \item Step-by-step execution display with action types
        \item Tool calls with expandable inputs/outputs
        \item Cypher query syntax highlighting
        \item Collapsible JSON output drawers
        \item Execution metrics (latency, token count)
    \end{itemize}
\end{itemize}

\subsection{Control Panel UI}
\label{subsec:control-panel-ui}

The Control Panel UI is a comprehensive Streamlit-based administration interface providing full API coverage with role-aware access control.

\paragraph{Technology Stack}
\begin{itemize}
    \item \textbf{Framework}: Streamlit (Python web apps)
    \item \textbf{Authentication}: Auth0 OAuth2/OIDC
    \item \textbf{HTTP Client}: httpx (async-capable)
    \item \textbf{Data Display}: Pandas DataFrames
    \item \textbf{Visualization}: Plotly (via Streamlit)
\end{itemize}

\paragraph{Application Structure}
\begin{lstlisting}[language=bash, caption={Control Panel directory structure.}]
ui_control_panel/
|-- app.py              # Main Streamlit entry point
|-- api.py              # API client with auth handling
|-- state.py            # Session state management
|-- components.py       # Reusable UI components
`-- views/
    |-- auth.py         # Authentication tab
    |-- dashboard.py    # Health monitoring
    |-- agents.py       # Agent runs & sessions
    |-- tools.py        # Tool discovery & invocation
    |-- models.py       # Model management
    |-- jobs.py         # Job management
    |-- tenants.py      # Multi-tenancy CRUD
    |-- admin.py        # Admin operations
    |-- cypher.py       # NL-to-Cypher workflow
    `-- explore.py      # API explorer
\end{lstlisting}

\paragraph{Authentication System}
The Control Panel supports four identity types with Auth0 integration:

\begin{table}[htbp]
    \centering
    \caption{Control Panel identity types.}
    \label{tab:control-panel-auth}
    \small
    \begin{tabular}{lll}
        \toprule
        \textbf{Identity Type} & \textbf{Auth Method} & \textbf{Use Case} \\
        \midrule
        Admin & Password Realm & Full platform administration \\
        User & Password Realm & Normal user operations \\
        Machine & Client Credentials & Service-to-service \\
        Custom & Manual token & Testing/debugging \\
        \bottomrule
    \end{tabular}
\end{table}

\paragraph{Scope-Based Access Control}
Features are enabled based on OAuth2 scopes in the active token:

\begin{table}[htbp]
    \centering
    \caption{Control Panel feature scopes.}
    \label{tab:control-panel-scopes}
    \small
    \begin{tabular}{ll}
        \toprule
        \textbf{Feature} & \textbf{Required Scopes} \\
        \midrule
        Dashboard & None (public health endpoints) \\
        Agents & \texttt{user:me} \\
        Tools (Safe) & \texttt{tools:invoke:basic} \\
        Tools (All) & \texttt{tools:invoke:all} \\
        Models (Create) & \texttt{admin:all} \\
        Tenants & \texttt{admin:all} \\
        Admin Operations & \texttt{admin:all} \\
        \bottomrule
    \end{tabular}
\end{table}

\paragraph{Tab-Based Navigation}
The Control Panel uses Streamlit's tab system to organize functionality:

\begin{enumerate}
    \item \textbf{Auth}: Role selection, token management, scope inspection
    \item \textbf{Dashboard}: System health, component status, KPI metrics
    \item \textbf{Agents}: Run creation, session management, step visualization
    \item \textbf{Tools}: Tool discovery, schema inspection, invocation testing
    \item \textbf{Models}: Provider management, instance configuration, defaults
    \item \textbf{Jobs}: Job creation, status monitoring, event streaming
    \item \textbf{Tenants}: Tenant CRUD, configuration management
    \item \textbf{Admin}: Database operations, system configuration
    \item \textbf{Cypher}: NL-to-Cypher testing, query execution
    \item \textbf{Explore}: OpenAPI specification, raw API requests
\end{enumerate}

\paragraph{API Client Architecture}
The API client provides robust backend communication:

\begin{lstlisting}[language=Python, caption={Control Panel API client pattern.}]
def normalize_endpoint(endpoint: str) -> str:
    """Ensure all endpoints start with /v1/"""
    if not endpoint.startswith('/'):
        endpoint = '/' + endpoint
    if not endpoint.startswith('/v1'):
        endpoint = '/v1' + endpoint
    return endpoint

def get_headers() -> dict:
    """Build request headers with auth and tenant context"""
    headers = {"Content-Type": "application/json"}
    token = get_active_token()
    if token:
        headers["Authorization"] = f"Bearer {token.value}"
    tenant_id = get_state().tenant_id
    if tenant_id:
        headers["X-Tenant-ID"] = tenant_id
    return headers
\end{lstlisting}

\paragraph{Dashboard Health Monitoring}
The dashboard tab monitors system health via multiple endpoints:

\begin{table}[htbp]
    \centering
    \caption{Dashboard health endpoints.}
    \label{tab:dashboard-health}
    \small
    \begin{tabular}{lll}
        \toprule
        \textbf{Endpoint} & \textbf{Purpose} & \textbf{Display} \\
        \midrule
        \texttt{/health/live} & Liveness probe & Status card \\
        \texttt{/health/ready} & Readiness probe & Status card \\
        \texttt{/health/startup} & Startup probe & Status card \\
        \texttt{/health/components} & Component health & Grid of cards \\
        \bottomrule
    \end{tabular}
\end{table}

\paragraph{Agent Execution Interface}
The agents tab provides a complete agent execution interface:
\begin{itemize}
    \item Run creator form with prompt, model selection, temperature, and max steps
    \item Session management for conversation continuity
    \item Real-time progress monitoring with step-by-step timeline
    \item Result display with final answer and execution metrics
    \item Polling with jitter to prevent thundering herd effects
\end{itemize}

\paragraph{NL-to-Cypher Testing}
The Cypher tab enables natural language to Cypher workflow testing:
\begin{enumerate}
    \item Natural language query input with example prompts
    \item Graph schema exploration and visualization
    \item Generated Cypher display with syntax highlighting
    \item Query execution with result rendering
    \item Query history for comparison and iteration
\end{enumerate}

\subsection{Dual-UI Architecture Analysis}
\label{subsec:dual-ui-analysis}

The decision to implement two separate user interfaces---rather than a single unified application---reflects a deliberate architectural choice based on persona analysis and technology fit.

\paragraph{Advantages of Dual-UI Approach}

\begin{description}
    \item[Persona Separation] The Chat UI serves end-users requiring conversational AI interaction, while the Control Panel serves administrators and operators requiring data exploration and system management. Each UI optimizes for its target persona without compromises.
    
    \item[Technology Specialization] Next.js excels at real-time, interactive experiences with sub-second feedback, while Streamlit excels at rapid development of data-driven dashboards. Using each framework for its strengths yields better user experiences than a single compromise solution.
    
    \item[Parallel Development] Separate codebases enable independent development by different team members or skill sets. Frontend specialists can focus on the Chat UI while backend developers contribute to the Control Panel.
    
    \item[Risk Isolation] Bugs or outages in one UI do not affect the other. The Chat UI can remain operational during Control Panel maintenance and vice versa.
    
    \item[Deployment Flexibility] Each UI can be deployed, scaled, and updated independently. The Chat UI can be served as static assets via CDN, while the Control Panel runs as a Python process.
\end{description}

\paragraph{Potential Limitations}

\begin{description}
    \item[Code Duplication] Both UIs implement their own API clients, authentication handling, and error formatting. Changes to API contracts require updates in both codebases.
    
    \item[State Inconsistency] Session state and user preferences are managed independently in each UI. A user switching between UIs may encounter inconsistent defaults or cached data.
    
    \item[Cognitive Overhead] Users who require both interfaces (e.g., developers testing agents and monitoring jobs) must context-switch between different visual languages and interaction patterns.
    
    \item[Testing Complexity] End-to-end tests must cover both interfaces, potentially with different testing frameworks (Playwright for Next.js, Selenium or direct Python tests for Streamlit).
    
    \item[Maintenance Burden] Feature additions (e.g., a new model field) may require changes in both UIs, increasing total development effort.
\end{description}

\paragraph{Mitigations Implemented}

\begin{itemize}
    \item \textbf{Shared Backend API}: The REST API serves as the single source of truth, ensuring both UIs reflect consistent data.
    \item \textbf{Consistent Model Selection}: Both UIs query the same \texttt{/v1/models/defaults} endpoint, ensuring default model consistency.
    \item \textbf{Common Health Endpoints}: Both UIs display health status from the same \texttt{/health/*} endpoints.
    \item \textbf{Documentation}: Clear documentation indicates when to use each UI based on task type.
\end{itemize}

\paragraph{Comparison Summary}
Table~\ref{tab:ui-comparison} summarizes the key differences between the two interfaces.

\begin{table}[htbp]
    \centering
    \caption{Comparison of Chat UI and Control Panel characteristics.}
    \label{tab:ui-comparison}
    \small
    \begin{tabularx}{\textwidth}{lXX}
        \toprule
        \textbf{Aspect} & \textbf{Agent Chat UI} & \textbf{Control Panel} \\
        \midrule
        Primary Persona & End users, researchers & Administrators, operators \\
        Interaction Style & Conversational, real-time & Dashboard, form-driven \\
        Technology & Next.js, React, TypeScript & Streamlit, Python \\
        State Management & Zustand (client-side) & Session state (server-side) \\
        Deployment & Static/SSR, CDN-cacheable & Python process, stateful \\
        Update Latency & Sub-second (polling) & Page refresh \\
        Development Speed & Moderate (typed, compiled) & Fast (Python, hot reload) \\
        \bottomrule
    \end{tabularx}
\end{table}

\subsection{Developer Experience and API Consumer Patterns}
\label{subsec:developer-experience}

Beyond the two provided UIs, the platform API is designed for consumption by external clients, scripts, and integrations. This subsection describes patterns for effective API usage.

\paragraph{RESTful Conventions}
The API follows standard REST conventions:
\begin{itemize}
    \item \textbf{Resource-Oriented URLs}: \texttt{/v1/agents}, \texttt{/v1/jobs/\{id\}}
    \item \textbf{HTTP Methods}: GET (read), POST (create), PUT/PATCH (update), DELETE (remove)
    \item \textbf{Status Codes}: 200 (success), 201 (created), 204 (no content), 4xx (client error), 5xx (server error)
    \item \textbf{Content Types}: \texttt{application/json} for requests/responses, \texttt{application/problem+json} for errors
\end{itemize}

\paragraph{Authentication Pattern}
API clients authenticate using JWT bearer tokens:

\begin{lstlisting}[language=bash, caption={API authentication example.}]
# Obtain token from Auth0
TOKEN=$(curl -s https://${AUTH0_DOMAIN}/oauth/token \
  -d "grant_type=client_credentials" \
  -d "client_id=${CLIENT_ID}" \
  -d "client_secret=${CLIENT_SECRET}" \
  -d "audience=${API_AUDIENCE}" | jq -r .access_token)

# Make authenticated request
curl -H "Authorization: Bearer ${TOKEN}" \
     -H "X-Tenant-ID: ${TENANT_ID}" \
     https://api.example.com/v1/agents
\end{lstlisting}

\paragraph{Error Handling}
Clients should handle RFC 7807 Problem Detail responses:

\begin{lstlisting}[language=Python, caption={Python client error handling.}]
import httpx

def create_agent_run(prompt: str, token: str) -> dict:
    response = httpx.post(
        f"{API_BASE}/v1/agent-runs",
        headers={"Authorization": f"Bearer {token}"},
        json={"prompt": prompt},
    )
    
    if response.status_code >= 400:
        problem = response.json()
        raise APIError(
            status=problem["status"],
            title=problem["title"],
            detail=problem.get("detail"),
        )
    
    return response.json()
\end{lstlisting}

\paragraph{Idempotent Requests}
For safe retries, clients should use idempotency keys:

\begin{lstlisting}[language=Python, caption={Idempotent request pattern.}]
import uuid

def create_job_idempotent(job_type: str, payload: dict, token: str) -> dict:
    idempotency_key = str(uuid.uuid4())
    
    response = httpx.post(
        f"{API_BASE}/v1/jobs",
        headers={
            "Authorization": f"Bearer {token}",
            "Idempotency-Key": idempotency_key,
        },
        json={"type": job_type, "payload": payload},
    )
    
    # Safe to retry with same idempotency_key
    return response.json()
\end{lstlisting}

\paragraph{Pagination Handling}
Clients should use cursor-based pagination for list endpoints:

\begin{lstlisting}[language=Python, caption={Cursor-based pagination example.}]
def list_all_runs(token: str) -> list[dict]:
    all_runs = []
    cursor = None
    
    while True:
        params = {"limit": 100}
        if cursor:
            params["cursor"] = cursor
        
        response = httpx.get(
            f"{API_BASE}/v1/agent-runs",
            headers={"Authorization": f"Bearer {token}"},
            params=params,
        )
        data = response.json()
        all_runs.extend(data["items"])
        
        if not data.get("has_more"):
            break
        cursor = data["next_cursor"]
    
    return all_runs
\end{lstlisting}

\paragraph{OpenAPI Documentation}
The API provides OpenAPI 3.1 documentation at \texttt{/openapi.json}:
\begin{itemize}
    \item Full endpoint documentation with request/response schemas
    \item Authentication requirements per endpoint
    \item Example requests and responses
    \item Interactive exploration via Swagger UI at \texttt{/docs}
\end{itemize}

\paragraph{SDK Generation}
The OpenAPI specification enables automatic SDK generation for various languages:

\begin{lstlisting}[language=bash, caption={SDK generation example.}]
# Generate Python client
openapi-generator-cli generate \
  -i https://api.example.com/openapi.json \
  -g python \
  -o ./sdk/python

# Generate TypeScript client
openapi-generator-cli generate \
  -i https://api.example.com/openapi.json \
  -g typescript-fetch \
  -o ./sdk/typescript
\end{lstlisting}
